
\documentclass[11pt,landscape]{article}
\usepackage[landscape,margin=1in]{geometry}
\usepackage{longtable,booktabs,array,tabularx,xcolor,hyperref,fancyhdr,graphicx,float}
\usepackage[T1]{fontenc}
\usepackage{helvet}
\renewcommand{\familydefault}{\sfdefault}

\definecolor{stridegreen}{HTML}{2E7D32}
\definecolor{strideblue}{HTML}{1565C0}
\definecolor{stridegray}{HTML}{757575}
\definecolor{csfbox}{HTML}{E3F2FD}
\definecolor{ferpabox}{HTML}{FFF3E0}
\definecolor{soxbox}{HTML}{F3E5F5}
\definecolor{gapred}{HTML}{C62828}

\hypersetup{colorlinks=true,linkcolor=strideblue,urlcolor=strideblue}

\pagestyle{fancy}
\fancyhf{}
\fancyhead[L]{\small\textcolor{stridegray}{Stride Security Reference Architecture V7}}
\fancyhead[R]{\small\textcolor{stridegray}{CONFIDENTIAL}}
\fancyfoot[C]{\thepage}
\fancyfoot[R]{\small March 2, 2026}
\renewcommand{\headrulewidth}{0.4pt}

\setlength{\parindent}{0pt}
\setlength{\parskip}{6pt}

\newcommand{\csfbox}[1]{\colorbox{csfbox}{\parbox{0.95\textwidth}{\small\textbf{NIST CSF 2.0 Alignment:} #1}}}
\newcommand{\ferpabox}[1]{\colorbox{ferpabox}{\parbox{0.95\textwidth}{\small\textbf{FERPA Relevance:} #1}}}
\newcommand{\soxbox}[1]{\colorbox{soxbox}{\parbox{0.95\textwidth}{\small\textbf{SOX §404 Relevance:} #1}}}

\begin{document}

% --- Title Page ---
\begin{titlepage}
\vspace*{1.5cm}
\begin{center}
{\Huge\bfseries\textcolor{stridegreen}{Stride Security Reference Architecture}}\\[0.5cm]
{\LARGE Version 7.0}\\[0.8cm]
{\large Enterprise Trust Security Reference Architecture}\\[0.3cm]
{\large Stride — K12 Inc.}\\[2cm]
{\normalsize Incorporating:}\\[0.2cm]
{\normalsize 22 D2 Architecture Diagrams (Landscape-Optimized)}\\
{\normalsize Per-Section NIST CSF 2.0 Cross-References}\\
{\normalsize FERPA and SOX §404 Regulatory Alignment}\\[2cm]
{\small Document Date: March 2, 2026}\\
{\small Classification: Confidential}\\[0.5cm]
{\small\textcolor{stridegray}{Supersedes: SRA V6 (February 27, 2026)}}\\
\end{center}
\end{titlepage}

\tableofcontents
\newpage


\section*{Diagram Legend and Conventions}
\addcontentsline{toc}{section}{Diagram Legend and Conventions}

\begin{center}
\includegraphics[width=0.85\textwidth]{/Volumes/Overflow/Repositories/d2.rescor.net/diagrams/security-reference-architecture/rendering_png_v2/00-legend-conventions.png}
\end{center}

\newpage
\section{Executive Summary}
\label{sra:s1}

\csfbox{GV.OC-01, GV.RM-01, GV.OV-01}

\vspace{0.3cm}

This Enterprise Trust Security Reference Architecture (TSRA) defines how Stride establishes, evaluates, enforces, degrades, and restores trust across identities, devices, workloads, APIs, data, and operations for cloud-native education technology platforms serving staff, students, and customers.

The architecture uses NIST CSF 2.0 as the primary outcomes framework and aligns to NIST SP 800-207 principles. It integrates current and planned security capabilities into a coherent model that is actionable, auditable, and resilient to tool changes, while remaining explicit about enforcement points and evidence sources.

\begin{itemize}

  \item Implement trust principles across all personas with identity-centric, risk-adaptive controls.

  \item Establish unified identity governance via SailPoint for provisioning, deprovisioning, access certification, and SoD.

  \item Protect privileged access using Delinea for vaulting, rotation, and session control; mature privileged role governance with Entra PIM (future).

  \item Maintain high-confidence device trust using CrowdStrike Falcon as the primary enforcement signal with Defender for Endpoint telemetry in monitor mode.

  \item Continuously assess exposure using Tenable VM for vulnerability management and Tenable.cs for CSPM posture and drift.

  \item Secure SDLC and supply-chain risk with Sonatype SCA, validate runtime security using Burp Suite DAST, and mature API security with Salt (future).

  \item Centralize operational trust and evidence via ServiceNow SecOps; mature CMDB and cloud tagging for ownership and criticality (future).

  \item Meet SOX and FERPA obligations with measurable controls and auditable evidence trails.

\end{itemize}


\begin{figure}[H]
\centering
\includegraphics[width=0.92\textwidth,height=0.72\textheight,keepaspectratio]{/Volumes/Overflow/Repositories/d2.rescor.net/diagrams/security-reference-architecture/rendering_png_v2/16-executive-summary-and-platform-stack.png}
\caption{16 Executive Summary And Platform Stack}
\end{figure}

\newpage

\section{Trust Security Architecture Overview}
\label{sra:s2}

\csfbox{GV.PO-01, PR.AA-01, PR.PS-01}

\vspace{0.3cm}

The TSRA implements the principle of 'never trust, always verify' as a continuous trust computation model. Trust is non-binary and time-bound: it is established with evidence, degraded when risk increases, and restored only after validation.

Zero Trust focuses on access decisioning. Trust Architecture extends this by governing the confidence required for specific actions over time, including data actions (view/export), privileged elevations, and API behaviors.


\begin{figure}[H]
\centering
\includegraphics[width=0.92\textwidth,height=0.72\textheight,keepaspectratio]{/Volumes/Overflow/Repositories/d2.rescor.net/diagrams/security-reference-architecture/rendering_png_v2/01-trust-security-architecture-overview.png}
\caption{01 Trust Security Architecture Overview}
\end{figure}


\begin{figure}[H]
\centering
\includegraphics[width=0.92\textwidth,height=0.72\textheight,keepaspectratio]{/Volumes/Overflow/Repositories/d2.rescor.net/diagrams/security-reference-architecture/rendering_png_v2/02-trust-logical-architecture-components.png}
\caption{02 Trust Logical Architecture Components}
\end{figure}

\newpage

\section{Personas and Trust Boundaries}
\label{sra:s3}

\csfbox{PR.AA-01, PR.AA-03, ID.AM-01}

\vspace{0.3cm}

\begin{itemize}

  \item Stride serves three primary user populations with differentiated identity, session requirements, and trust assumptions.

  \item Staff are the highest-risk population due to access to sensitive data, administrative systems, and infrastructure.

  \item Students represent the largest population. Access is governed through Okta CIAM with lifecycle driven by SIS via SailPoint. Trust is primarily contextual and behavioral due to unmanaged endpoints.

  \item Customers (parents/guardians) use self-service registration with verification. Step-up trust is required for payments and PII access.

\end{itemize}


\begin{figure}[H]
\centering
\includegraphics[width=0.92\textwidth,height=0.72\textheight,keepaspectratio]{/Volumes/Overflow/Repositories/d2.rescor.net/diagrams/security-reference-architecture/rendering_png_v2/03-persona-trust-boundaries.png}
\caption{03 Persona Trust Boundaries}
\end{figure}

\newpage

\section{Identity and Access Architecture}
\label{sra:s4}

\csfbox{PR.AA-01, PR.AA-02, PR.PS-01}

\vspace{0.3cm}

\ferpabox{FERPA §99.31 (consent exceptions for legitimate educational interest) and §99.32 (disclosure records) require identity controls that enforce role-based access aligned with educational purpose.}

\vspace{0.2cm}

\soxbox{SOX §404 AC-1 through AC-5: logical access provisioning, periodic reviews, PAM, MFA, and service account controls directly implement IAM architecture requirements.}

\vspace{0.2cm}

\vspace{0.2cm}

\colorbox{gapred!10}{\parbox{0.95\textwidth}{\small\textbf{Gap ---} FERPA Consent and Disclosure Controls: §99.30 (consent for disclosures) and §99.35 (access by state/local officials) lack explicit architectural controls for consent management and disclosure tracking beyond role-based access. NIST CSF PR.AA-04 (identity assertions managed) and PR.AA-06 (access rights managed) are not addressed in the current SRA subcategory mapping. A formal consent-management workflow and disclosure-audit mechanism should be established within the SailPoint/Entra governance layer.}}

\vspace{0.2cm}

The identity architecture is hybrid: Entra ID is authoritative for workforce identities and Okta CIAM is authoritative for external identities. SailPoint provides unified governance and attestation across both.


\begin{figure}[H]
\centering
\includegraphics[width=0.92\textwidth,height=0.72\textheight,keepaspectratio]{/Volumes/Overflow/Repositories/d2.rescor.net/diagrams/security-reference-architecture/rendering_png_v2/04-staff-authentication-flow.png}
\caption{04 Staff Authentication Flow}
\end{figure}


\begin{figure}[H]
\centering
\includegraphics[width=0.92\textwidth,height=0.72\textheight,keepaspectratio]{/Volumes/Overflow/Repositories/d2.rescor.net/diagrams/security-reference-architecture/rendering_png_v2/05-staff-conditional-access-policy-matrix.png}
\caption{05 Staff Conditional Access Policy Matrix}
\end{figure}


\begin{figure}[H]
\centering
\includegraphics[width=0.92\textwidth,height=0.72\textheight,keepaspectratio]{/Volumes/Overflow/Repositories/d2.rescor.net/diagrams/security-reference-architecture/rendering_png_v2/06-student-authentication-flow.png}
\caption{06 Student Authentication Flow}
\end{figure}


\begin{figure}[H]
\centering
\includegraphics[width=0.92\textwidth,height=0.72\textheight,keepaspectratio]{/Volumes/Overflow/Repositories/d2.rescor.net/diagrams/security-reference-architecture/rendering_png_v2/07-customer-authentication-flow.png}
\caption{07 Customer Authentication Flow}
\end{figure}


\begin{figure}[H]
\centering
\includegraphics[width=0.92\textwidth,height=0.72\textheight,keepaspectratio]{/Volumes/Overflow/Repositories/d2.rescor.net/diagrams/security-reference-architecture/rendering_png_v2/08-identity-and-access-architecture.png}
\caption{08 Identity And Access Architecture}
\end{figure}

\newpage

\section{Privileged Trust (Delinea and Entra PIM)}
\label{sra:s5}

\csfbox{PR.AA-05, PR.PS-02, DE.CM-03}

\vspace{0.3cm}

\soxbox{SOX §404 AC-3: privileged access management with JIT elevation and session recording.}

\vspace{0.2cm}

Privileged access is controlled through Delinea for credential vaulting, rotation, and session control for privileged accounts, with Entra PIM planned to mature privileged role governance for Entra roles.


\begin{figure}[H]
\centering
\includegraphics[width=0.92\textwidth,height=0.72\textheight,keepaspectratio]{/Volumes/Overflow/Repositories/d2.rescor.net/diagrams/security-reference-architecture/rendering_png_v2/09-privileged-trust-architecture.png}
\caption{09 Privileged Trust Architecture}
\end{figure}

\newpage

\section{Device Trust (CrowdStrike and Defender)}
\label{sra:s6}

\csfbox{PR.PS-03, DE.CM-01}

\vspace{0.3cm}

\begin{itemize}

  \item CrowdStrike Falcon provides primary endpoint protection and device trust signals. Microsoft Defender for Endpoint is deployed in monitor mode to provide supplemental telemetry.

\end{itemize}


\begin{figure}[H]
\centering
\includegraphics[width=0.92\textwidth,height=0.72\textheight,keepaspectratio]{/Volumes/Overflow/Repositories/d2.rescor.net/diagrams/security-reference-architecture/rendering_png_v2/10-device-trust-architecture.png}
\caption{10 Device Trust Architecture}
\end{figure}

\newpage

\section{Exposure and Posture (Tenable VM and CSPM)}
\label{sra:s7}

\csfbox{ID.RA-01, DE.CM-02, PR.PS-04}

\vspace{0.3cm}

\begin{itemize}

  \item Tenable VM provides vulnerability identification across infrastructure and endpoints. Tenable.cs provides CSPM posture coverage and drift detection.

\end{itemize}


\begin{figure}[H]
\centering
\includegraphics[width=0.92\textwidth,height=0.72\textheight,keepaspectratio]{/Volumes/Overflow/Repositories/d2.rescor.net/diagrams/security-reference-architecture/rendering_png_v2/11-exposure-and-posture-architecture.png}
\caption{11 Exposure And Posture Architecture}
\end{figure}

\newpage

\section{Application and API Trust (Sonatype, Burp Suite, Salt)}
\label{sra:s8}

\csfbox{PR.DS-01, PR.PS-05, DE.CM-09}

\vspace{0.3cm}

\soxbox{SOX §404 PD-1 through PD-2 and CM-2 through CM-6: SDLC methodology, security testing in CI/CD, environment segregation, and change approval.}

\vspace{0.2cm}

\vspace{0.2cm}

\colorbox{gapred!10}{\parbox{0.95\textwidth}{\small\textbf{Gap ---} API Authentication Controls: The current architecture addresses application-layer testing (SAST/DAST/SCA) and future API behavioral analytics (Salt) but does not prescribe specific API authentication standards (e.g., OAuth 2.0/OIDC scoped tokens vs. static API keys). A formal API authentication standard should be developed under ISP 4.1 (PR.AA) to define acceptable mechanisms, token lifecycle management, and mutual-TLS requirements for service-to-service communication. The ASR Questionnaire (V4) now includes an API authentication assessment question to capture the current state per application. Additionally, SOX CM-1 (change management policy documentation) is referenced but the architecture does not prescribe a formal change management policy artefact. PD-3 (security testing sign-off) and PD-4 (system conversion testing) lack documented verification gates for security test acceptance and data-migration validation.}}

\vspace{0.2cm}

\begin{itemize}

  \item Sonatype SCA manages dependency risk; Burp Suite DAST validates runtime security; Salt (future) provides API discovery and behavior analytics.

\end{itemize}


\begin{figure}[H]
\centering
\includegraphics[width=0.92\textwidth,height=0.72\textheight,keepaspectratio]{/Volumes/Overflow/Repositories/d2.rescor.net/diagrams/security-reference-architecture/rendering_png_v2/12-application-and-api-trust-architecture.png}
\caption{12 Application And Api Trust Architecture}
\end{figure}

\newpage

\section{Operational Trust (ServiceNow SecOps, CMDB, Tagging)}
\label{sra:s9}

\csfbox{DE.CM-06, RS.MA-01, RC.RP-01}

\vspace{0.3cm}

\soxbox{SOX §404 CO-1 through CO-4: job scheduling, backup/recovery, incident management, and batch processing integrity.}

\vspace{0.2cm}

\vspace{0.2cm}

\colorbox{gapred!10}{\parbox{0.95\textwidth}{\small\textbf{Gap ---} Operational Monitoring and Event Analysis: DE.CM-06 (external service provider activities monitored) is In Progress---ServiceNow integration with third-party telemetry feeds is not yet operational. SOX CO-1 (job scheduling) and CO-4 (batch processing integrity) lack specific ServiceNow automation evidence. The NIST CSF DE.AE (Adverse Event Analysis) function---covering DE.AE-02 through DE.AE-08---is entirely absent from the SRA subcategory mapping; correlation, enrichment, and escalation of security events require formal architectural treatment.}}

\vspace{0.2cm}

\begin{itemize}

  \item ServiceNow SecOps is the central record for security incidents, vulnerability response, and evidence trails. CMDB and tagging enhance ownership and prioritization (future).

\end{itemize}


\begin{figure}[H]
\centering
\includegraphics[width=0.92\textwidth,height=0.72\textheight,keepaspectratio]{/Volumes/Overflow/Repositories/d2.rescor.net/diagrams/security-reference-architecture/rendering_png_v2/13-operational-trust-architecture.png}
\caption{13 Operational Trust Architecture}
\end{figure}

\newpage

\section{NIST CSF 2.0 Outcome Mapping}
\label{sra:s10}

\csfbox{GV.OV-01, ID.RA-01, PR.AA-01, DE.CM-01, RS.MA-01, RC.RP-01}

\vspace{0.3cm}

\begin{itemize}

  \item NIST CSF 2.0 is used as the primary outcomes framework across Govern, Identify, Protect, Detect, Respond, and Recover.

\end{itemize}


\begin{figure}[H]
\centering
\includegraphics[width=0.92\textwidth,height=0.72\textheight,keepaspectratio]{/Volumes/Overflow/Repositories/d2.rescor.net/diagrams/security-reference-architecture/rendering_png_v2/14-nist-csf-outcome-mapping.png}
\caption{14 Nist Csf Outcome Mapping}
\end{figure}

\newpage

\section{Implementation Roadmap}
\label{sra:s11}

\csfbox{GV.SC-01, GV.RM-02}

\vspace{0.3cm}

\vspace{0.2cm}

\colorbox{gapred!10}{\parbox{0.95\textwidth}{\small\textbf{Gap ---} Supply Chain and Risk Appetite: GV.SC-01 (supply chain risk management program) is Planned and not yet initiated. GV.RM-02 (risk appetite and tolerance statements) is In Progress. Additional CSF subcategories GV.SC-02 through GV.SC-10---covering supplier due-diligence, contractual security requirements, supply chain integrity verification, and post-acquisition monitoring---are not mapped in the SRA. A formal third-party risk management program aligned to ISP 3.2 and NIST SP 800-161r1 should be established.}}

\vspace{0.2cm}

\begin{itemize}

  \item Roadmap phases describe maturity sequencing across trust domains.

\end{itemize}


\begin{figure}[H]
\centering
\includegraphics[width=0.92\textwidth,height=0.72\textheight,keepaspectratio]{/Volumes/Overflow/Repositories/d2.rescor.net/diagrams/security-reference-architecture/rendering_png_v2/15-implementation-roadmap.png}
\caption{15 Implementation Roadmap}
\end{figure}

\newpage

\section{Key Performance Indicators (KPIs)}
\label{sra:s12}

\csfbox{GV.OV-02, DE.CM-09, RS.MA-04}

\vspace{0.3cm}

\begin{itemize}

  \item KPIs track adoption and effectiveness of trust controls (examples: MFA adoption, vuln SLAs, MTTD/MTTR).

\end{itemize}


\begin{figure}[H]
\centering
\includegraphics[width=0.92\textwidth,height=0.72\textheight,keepaspectratio]{/Volumes/Overflow/Repositories/d2.rescor.net/diagrams/security-reference-architecture/rendering_png_v2/17-kpi-framework.png}
\caption{17 Kpi Framework}
\end{figure}

\newpage

\section{Regulatory Compliance Alignment (SOX, FERPA)}
\label{sra:s13}

\csfbox{GV.RM-03, GV.OV-03}

\vspace{0.3cm}

\ferpabox{This section provides the comprehensive FERPA alignment matrix covering all applicable sections of 34 CFR Part 99.}

\vspace{0.2cm}

\soxbox{This section provides the comprehensive SOX §404 ITGC alignment matrix.}

\vspace{0.2cm}

\vspace{0.2cm}

\colorbox{gapred!10}{\parbox{0.95\textwidth}{\small\textbf{Gap ---} FERPA Coverage Gaps and Awareness Training: FERPA 34 CFR §§99.4 (student rights notification), 99.7 (annual FERPA notice requirements), 99.12 (limitations on re-disclosure), 99.33 (conditions on re-disclosure to third parties), and 99.37 (directory information opt-out procedures) lack specific architectural controls or evidence sources in the SRA. NIST CSF PR.AT (Awareness and Training) subcategories PR.AT-01 and PR.AT-02 are not addressed in the operational control architecture; a formal security awareness and role-based training program should be mapped.}}

\vspace{0.2cm}

\begin{itemize}

  \item SOX: Access governance, change hygiene, and auditable evidence. FERPA: Protect education records via least privilege, governance, and auditable incident handling.

\end{itemize}


\begin{figure}[H]
\centering
\includegraphics[width=0.92\textwidth,height=0.72\textheight,keepaspectratio]{/Volumes/Overflow/Repositories/d2.rescor.net/diagrams/security-reference-architecture/rendering_png_v2/18-regulatory-compliance-alignment.png}
\caption{18 Regulatory Compliance Alignment}
\end{figure}

\newpage


\appendix
\section{APPENDIX A --- Acronyms and Definitions}
\label{app:acronyms}

\begin{center}
\includegraphics[width=0.92\textwidth,height=0.72\textheight,keepaspectratio]{/Volumes/Overflow/Repositories/d2.rescor.net/diagrams/security-reference-architecture/rendering_png_v2/19-acronyms-and-definitions-map.png}
\end{center}

\vspace{0.5cm}

\begin{longtable}{lp{7in}}
\toprule
\textbf{Term} & \textbf{Definition} \\
\midrule
\endhead
CA & Conditional Access --- policy-based access control in Microsoft Entra that evaluates signals (user, device, location, risk) to enforce authentication requirements. \\
CIAM & Customer Identity and Access Management --- IAM specifically for external/customer-facing authentication (e.g., Okta CIC for students and families). \\
CMDB & Configuration Management Database --- central repository of IT asset records and relationships (ServiceNow CMDB). \\
CSF & Cybersecurity Framework --- NIST framework (v2.0) organizing cybersecurity outcomes across Govern, Identify, Protect, Detect, Respond, and Recover functions. \\
CSPM & Cloud Security Posture Management --- continuous monitoring of cloud configurations against security baselines (Tenable Cloud Security). \\
EDR & Endpoint Detection and Response --- real-time endpoint monitoring and threat response (CrowdStrike Falcon). \\
FERPA & Family Educational Rights and Privacy Act --- 34 CFR Part 99, protecting student education records. \\
FIDO2 & Fast Identity Online 2 --- passwordless authentication standard using hardware security keys or platform authenticators. \\
IAM & Identity and Access Management --- policies and technologies ensuring the right individuals access the right resources at the right time. \\
IGA & Identity Governance and Administration --- lifecycle management, access certification, and compliance reporting for identities. \\
ITGC & IT General Controls --- SOX \S 404 controls over access, change management, computer operations, and program development. \\
JEA & Just-Enough Administration --- restricting privileged sessions to the minimum set of commands required for the task. \\
JIT & Just-in-Time Access --- time-bound privilege elevation that expires automatically (Entra PIM, Delinea). \\
MDM & Mobile Device Management --- device enrollment, policy enforcement, and remote wipe capabilities. \\
MFA & Multi-Factor Authentication --- requiring two or more verification factors for access. \\
MTTR & Mean Time to Remediate/Respond --- average time between detection of a vulnerability or incident and its resolution. \\
NIST & National Institute of Standards and Technology --- U.S. federal agency publishing cybersecurity standards and guidelines. \\
PAM & Privileged Access Management --- controls for discovering, managing, and auditing privileged accounts and sessions (Delinea). \\
PIM & Privileged Identity Management --- role-based JIT elevation with approval workflows (Microsoft Entra PIM). \\
RBAC & Role-Based Access Control --- access permissions assigned by organizational role rat her than individual identity. \\
SecOps & Security Operations --- the team and processes for monitoring, detecting, and responding to security events. \\
SLA & Service Level Agreement --- formal commitment on service availability, response time, or remediation timelines. \\
SOX & Sarbanes-Oxley Act --- U.S. legislation requiring internal controls over financial reporting (\S 404). \\
SRA & Security Reference Architecture --- this document; the enterprise trust and security architecture blueprint. \\
TOGAF & The Open Group Architecture Framework --- enterprise architecture methodology used for platform and security architecture alignment. \\
VM & Vulnerability Management --- continuous process of identifying, classifying, prioritizing, and remediating vulnerabilities (Tenable). \\
WHfB & Windows Hello for Business --- passwordless, certificate-based authentication using biometrics or PIN bound to a device. \\
ZT & Zero Trust --- security model requiring continuous verification of every user, device, and transaction regardless of network location (NIST SP 800-207). \\
\bottomrule
\end{longtable}

\newpage

\section{APPENDIX B --- Document Control}
\label{app:doccontrol}

\begin{center}
\includegraphics[width=0.92\textwidth,height=0.72\textheight,keepaspectratio]{/Volumes/Overflow/Repositories/d2.rescor.net/diagrams/security-reference-architecture/rendering_png_v2/20-document-control-metadata.png}
\end{center}

\vspace{0.5cm}

\begin{tabularx}{\textwidth}{lX}
\toprule
\textbf{Attribute} & \textbf{Value} \\
\midrule
Document Title & Stride Security Reference Architecture \\
Version & 7.0 \\
Date & March 3, 2026 \\
Classification & Confidential \\
Author & Security Architecture / GRC \\
Supersedes & SRA V6 (February 27, 2026) \\
\bottomrule
\end{tabularx}

\vspace{0.5cm}

\begin{tabularx}{\textwidth}{llX}
\toprule
\textbf{Version} & \textbf{Date} & \textbf{Change Description} \\
\midrule
V6.0 & February 27, 2026 & Initial trust security architecture; 22 D2 diagrams; NIST CSF 2.0 alignment. \\
V7.0 & March 3, 2026 & Enhanced appendices (A--E); per-section FERPA/SOX alignment; gap annotations; consolidated gap register; diagram routing improvements; acronym and definitions table; NIST CSF and evidence traceability appendices. \\
\bottomrule
\end{tabularx}

\newpage

\section{APPENDIX C --- NIST CSF 2.0 Mapping and Implementation Status}
\label{app:csf_mapping}

This appendix maps NIST Cybersecurity Framework 2.0 subcategories referenced in this SRA
to the corresponding architecture sections and their current implementation status.

\vspace{0.3cm}

\begin{longtable}{lp{3in}p{2in}l}
\toprule
\textbf{CSF Subcategory} & \textbf{Outcome Description} & \textbf{SRA Section} & \textbf{Status} \\
\midrule
\endhead
GV.OC-01 & Organizational context for cybersecurity risk management & §1 Executive Summary & \textcolor{stridegreen}{Implemented} \\
GV.RM-01 & Risk management objectives established and communicated & §1 Executive Summary & \textcolor{stridegreen}{Implemented} \\
GV.RM-02 & Risk appetite and tolerance statements defined & §11 Implementation Roadmap & \textcolor{orange}{In Progress} \\
GV.RM-03 & Risk management activities integrated across enterprise & §13 Regulatory Compliance & \textcolor{stridegreen}{Implemented} \\
GV.OV-01 & Cybersecurity strategy outcomes assessed & §1, §10 NIST CSF Mapping & \textcolor{stridegreen}{Implemented} \\
GV.OV-02 & Cybersecurity performance evaluated against KPIs & §12 KPIs & \textcolor{stridegreen}{Implemented} \\
GV.OV-03 & Regulatory requirements integrated into governance & §13 Regulatory Compliance & \textcolor{stridegreen}{Implemented} \\
GV.PO-01 & Cybersecurity policy established and communicated & §2 Architecture Overview & \textcolor{stridegreen}{Implemented} \\
GV.SC-01 & Supply chain risk management program established & §11 Implementation Roadmap & \textcolor{gapred}{Planned} \\
ID.AM-01 & Asset inventories maintained & §3 Personas and Trust Boundaries & \textcolor{stridegreen}{Implemented} \\
ID.RA-01 & Vulnerabilities identified, validated, and recorded & §7 Exposure and Posture, §10 & \textcolor{stridegreen}{Implemented} \\
PR.AA-01 & Identities and credentials managed & §4 Identity and Access & \textcolor{stridegreen}{Implemented} \\
PR.AA-02 & Identity proofing and binding to credentials & §4 Identity and Access & \textcolor{stridegreen}{Implemented} \\
PR.AA-03 & Users, services, and hardware authenticated & §3, §4 Identity and Access & \textcolor{stridegreen}{Implemented} \\
PR.AA-05 & Least privilege and separation of duties enforced & §5 Privileged Trust & \textcolor{stridegreen}{Implemented} \\
PR.DS-01 & Data-at-rest and data-in-transit protected & §8 Application and API Trust & \textcolor{stridegreen}{Implemented} \\
PR.PS-01 & Configuration baselines maintained & §2, §4 Architecture Overview & \textcolor{stridegreen}{Implemented} \\
PR.PS-02 & Software maintained and replaced per policy & §5 Privileged Trust & \textcolor{stridegreen}{Implemented} \\
PR.PS-03 & Hardware maintained and replaced per policy & §6 Device Trust & \textcolor{stridegreen}{Implemented} \\
PR.PS-04 & Log integrity and availability ensured & §7 Exposure and Posture & \textcolor{stridegreen}{Implemented} \\
PR.PS-05 & Installation and execution of unauthorized software prevented & §8 Application and API & \textcolor{orange}{In Progress} \\
DE.CM-01 & Networks monitored for anomalies and attacks & §6 Device Trust & \textcolor{stridegreen}{Implemented} \\
DE.CM-02 & Physical environment monitored for anomalies & §7 Exposure and Posture & \textcolor{stridegreen}{Implemented} \\
DE.CM-03 & Personnel activity monitored for anomalies & §5 Privileged Trust & \textcolor{stridegreen}{Implemented} \\
DE.CM-06 & External service provider activities monitored & §9 Operational Trust & \textcolor{orange}{In Progress} \\
DE.CM-09 & Computing hardware and software monitored & §8, §12 KPIs & \textcolor{stridegreen}{Implemented} \\
RS.MA-01 & Incident response plan executed & §9 Operational Trust & \textcolor{stridegreen}{Implemented} \\
RS.MA-04 & Incident response performance evaluated & §12 KPIs & \textcolor{stridegreen}{Implemented} \\
RC.RP-01 & Recovery plan executed during/after incident & §9 Operational Trust & \textcolor{stridegreen}{Implemented} \\
\bottomrule
\end{longtable}

{\small\textit{Status definitions: \textcolor{stridegreen}{Implemented} = controls active and evidenced;
\textcolor{orange}{In Progress} = controls partially deployed or in current roadmap phase;
\textcolor{gapred}{Planned} = controls identified but not yet initiated.}}

\newpage

\section{APPENDIX D --- Evidence Traceability Matrix}
\label{app:evidence}

This appendix identifies the primary evidence artifacts for each NIST CSF 2.0 subcategory,
supporting audit readiness and attestation preparation.

\vspace{0.3cm}

\begin{longtable}{lp{3.5in}p{2in}}
\toprule
\textbf{CSF Subcategory} & \textbf{Evidence Artifacts} & \textbf{Evidence Owner} \\
\midrule
\endhead
GV.OC-01 & SRA document, governance charter & Security Architecture \\
GV.RM-01 & Risk register, risk appetite statement & GRC / Risk Management \\
GV.OV-01 & CSF outcome mapping (§10), annual review records & Security Architecture \\
GV.OV-02 & KPI dashboard (§12), quarterly metrics reports & Security Operations \\
GV.PO-01 & ISP 1.0–12.0 policy library, IISP 2.0 & Information Security \\
ID.AM-01 & ServiceNow CMDB asset inventory, Entra directory & IT Operations \\
ID.RA-01 & Tenable scan reports, risk assessment records & Vulnerability Management \\
PR.AA-01 & Entra ID configurations, Okta CIC tenant configs & Identity Engineering \\
PR.AA-02 & FIDO2 key enrollment records, MFA registration logs & Identity Engineering \\
PR.AA-03 & Conditional Access policy exports, sign-in logs & Identity Engineering \\
PR.AA-05 & Delinea session recordings, PIM activation logs & Privileged Access Mgmt \\
PR.DS-01 & TLS configurations, encryption-at-rest settings & Application Security \\
PR.PS-01 & CIS benchmark scan results, baseline configs & Infrastructure Security \\
PR.PS-03 & CrowdStrike sensor deployment reports, Defender status & Endpoint Security \\
PR.PS-05 & Sonatype policy results, application allowlists & Application Security \\
DE.CM-01 & CrowdStrike detection logs, network alert reports & Security Operations \\
DE.CM-03 & PIM audit logs, privileged session recordings & Security Operations \\
DE.CM-09 & Burp Suite scan reports, Salt Security API findings & Application Security \\
RS.MA-01 & ServiceNow incident records, runbook executions & Incident Response \\
RS.MA-04 & MTTR reports, incident post-mortem records & Incident Response \\
RC.RP-01 & Recovery test results, backup verification logs & Business Continuity \\
\bottomrule
\end{longtable}

{\small\textit{Evidence sources should be aligned to the formal control library and refreshed
during each assessment cycle. Artifacts are referenced in the control mapping tables
throughout this SRA.}}

\newpage

\section{APPENDIX E --- Consolidated Gap Register}
\label{app:gaps}

This appendix consolidates all identified gaps across NIST CSF 2.0 subcategory coverage,
FERPA regulatory requirements, and SOX §404 ITGC controls. Each gap references the SRA
section where the deficiency is annotated and includes a priority for remediation planning.

\vspace{0.3cm}

\begin{longtable}{lp{1.2in}lp{3.5in}l}
\toprule
\textbf{Gap ID} & \textbf{CSF / Regulatory Ref} & \textbf{SRA §} & \textbf{Description} & \textbf{Priority} \\
\midrule
\endhead
GAP-01 & PR.AA-04, PR.AA-06 & §4 & FERPA §99.30/§99.35 consent and disclosure controls lack explicit architectural enforcement; PR.AA-04 and PR.AA-06 unmapped. & \textcolor{gapred}{High} \\
GAP-02 & PR.AA (API) & §8 & No formal API authentication standard (OAuth 2.0/OIDC vs. static keys); mutual-TLS and token lifecycle undefined. & \textcolor{gapred}{High} \\
GAP-03 & CM-1, PD-3, PD-4 & §8 & SOX CM-1 change management policy artefact not prescribed; PD-3/PD-4 security testing sign-off and conversion validation gates missing. & \textcolor{orange}{Medium} \\
GAP-04 & DE.CM-06 & §9 & External service provider monitoring In Progress — ServiceNow integration with third-party telemetry not yet operational. & \textcolor{gapred}{High} \\
GAP-05 & CO-1, CO-4 & §9 & SOX CO-1 (job scheduling) and CO-4 (batch processing integrity) lack specific ServiceNow automation evidence. & \textcolor{orange}{Medium} \\
GAP-06 & DE.AE-02–08 & §9 & Entire DE.AE (Adverse Event Analysis) function absent from SRA mapping; event correlation, enrichment, and escalation unaddressed. & \textcolor{gapred}{High} \\
GAP-07 & GV.SC-01 & §11 & Supply chain risk management program Planned but not initiated. & \textcolor{gapred}{High} \\
GAP-08 & GV.RM-02 & §11 & Risk appetite and tolerance statements In Progress. & \textcolor{orange}{Medium} \\
GAP-09 & GV.SC-02–10 & §11 & Supplier due-diligence, contractual security, supply chain integrity, and post-acquisition monitoring subcategories unmapped. & \textcolor{gapred}{High} \\
GAP-10 & §99.4, §99.7, §99.12, §99.33, §99.37 & §13 & Five FERPA sections lack specific architectural controls or evidence sources in the SRA. & \textcolor{gapred}{High} \\
GAP-11 & PR.AT-01, PR.AT-02 & §13 & Awareness and Training subcategories not addressed in operational control architecture. & \textcolor{orange}{Medium} \\
GAP-12 & PR.PS-05 & §8 & Unauthorized software prevention In Progress — Sonatype policy enforcement not fully deployed. & \textcolor{orange}{Medium} \\
GAP-13 & GV.RM, PR.AT, DE.AE & ASR & ASR questionnaire lacks questions covering risk management governance, awareness training, and adverse event analysis domains. & \textcolor{gapred}{High} \\
\bottomrule
\end{longtable}

{\small\textit{Gaps are derived from the comprehensive control-target analysis performed in
March 2026 comparing NIST CSF 2.0 (100 subcategories), FERPA 34 CFR Part 99 (18 sections),
and SOX §404 ITGC (19 controls) against the SRA and ASR questionnaire coverage.
Priority reflects the assessed impact on audit readiness and regulatory compliance.}}

\end{document}
